% =====================================================================
\section{Función 91 de McCarthy, inducción modificada y Ackermann}
% =====================================================================

\begin{prob}
Con $f(x)=x-10$ si $x>100$ y $f(x)=f(f(x+11))$ si $x\le100$:
(a) desarrolla a mano $f(99)$ y $f(95)$ y cuenta cuántas llamadas se hacen en cada caso;
(b) desarrolla $f(89)$ mostrando solo los pasos intermedios importantes;
(c) calcula $f(105)$ y $f(101)$.
\end{prob}
\begin{sol}
(a) $f(99)=f(f(110))=f(100)=f(f(111))=f(101)=91$. Llamadas: $f(99),f(110),f(100),f(111),f(101)$: \textbf{5}.\\
$f(95)=f(f(106))=f(96)=f(f(107))=f(97)=\dots=f(100)=f(f(111))=f(101)=91$. Se llama a $f(95),\dots,f(101)$ (7) y a $f(106),\dots,f(111)$ (6): \textbf{13} llamadas.

(b) $f(89)=f(f(100))$; como $f(100)=91$ (3 llamadas), queda $f(91)=f(f(102))=f(92)=\dots=f(100)=f(101)=91$. (Total: 25 llamadas.)

(c) $f(105)=95$; \ $f(101)=91$. En general $f(x)=x-10$ para $x>100$ y $f(x)=91$ para $x\le101$.
\end{sol}

\begin{prob}
Enuncia el principio de inducción matemática modificado y úsalo para demostrar que $f(x)=91$ para todo entero $x\le100$. Explica por qué la inducción simple (``si vale en $n$ vale en $n+1$'') no basta.
\end{prob}
\begin{sol}
\textbf{Principio modificado.} Sea $k$ un entero fijo. Si (1) $P(k)$ es verdadera y (2) para todo $n<k$, [$P(m)$ verdadera para todo $m$ con $n<m\le k$] implica $P(n)$, entonces $P(n)$ vale para todo $n\le k$.

\textbf{Demostración} con $k=100$ y $P(x): f(x)=91$.

\emph{Paso básico:} $f(100)=f(f(111))=f(101)=91$ (pues $111>100$ y $101>100$).

\emph{Hipótesis:} sea $x<100$ y supongamos $f(m)=91$ para todo $m$ con $x<m\le100$.

\emph{Paso inductivo:} como $x\le100$, $f(x)=f(f(x+11))$.
\begin{itemize}
\item \textbf{Caso a)} $x+11>100$ (es decir, $90\le x\le99$): $f(x+11)=x+1$, luego $f(x)=f(x+1)$. Como $x<x+1\le100$, por H.I.\ $f(x+1)=91$.
\item \textbf{Caso b)} $x+11\le100$ ($x\le89$): como $x<x+11\le100$, por H.I.\ $f(x+11)=91$, luego $f(x)=f(91)$. Como $x\le89<91\le100$, por H.I.\ $f(91)=91$.
\end{itemize}
En ambos casos $f(x)=91$. Por el principio modificado, $f(x)=91$ para todo $x\le100$. $\blacksquare$

\textbf{Por qué no basta la inducción simple:} el caso b) necesita el valor de $f$ en \emph{dos} puntos mayores que $x$ ($x+11$ y $91$), no solo en $x+1$; la hipótesis debe cubrir todo el tramo $(x,100]$. Además la recursión ``se aleja'' (el argumento sube) antes de resolverse, así que no hay un contador que baje de uno en uno.
\end{sol}

\begin{prob}[\simil]
Sea $g(x)=x-2$ si $x>20$ y $g(x)=g(g(x+3))$ si $x\le20$. Calcula $g(20)$, $g(19)$, $g(15)$; conjetura el valor de $g(x)$ para $x\le20$ y demuéstralo con inducción modificada.
\end{prob}
\begin{sol}
$g(20)=g(g(23))=g(21)=19$; \ $g(19)=g(g(22))=g(20)=19$; \ $g(15)=g(g(18))=g(19)=19$ (usando $g(18)=g(g(21))=g(19)=19$).
\textbf{Conjetura:} $g(x)=19$ para todo $x\le20$.

\textbf{Paso básico} ($k=20$): $g(20)=19$ (calculado).
\textbf{H.I.:} para $x<20$, $g(m)=19$ para todo $m$ con $x<m\le20$.
\textbf{Paso:} $g(x)=g(g(x+3))$.
\begin{itemize}
\item Caso a) $x+3>20$ ($x\in\{18,19\}$): $g(x+3)=x+1$, así $g(x)=g(x+1)$ con $x<x+1\le20$: por H.I.\ vale 19.
\item Caso b) $x+3\le20$ ($x\le17$): $x<x+3\le20\Rightarrow g(x+3)=19$; luego $g(x)=g(19)$ y $x\le17<19\le20\Rightarrow g(19)=19$.
\end{itemize}
Por inducción modificada, $g(x)=19$ para todo $x\le20$. $\blacksquare$ \ (Patrón general del tipo McCarthy: techo $c$, resta $a$, suma $b>a$: valor $c+b-2a$; aquí $20+3-4=19$ y $100+11-20=91$.)
\end{sol}

\begin{prob}
Con la definición $A(0,y)=y+1$, $A(x+1,0)=A(x,1)$, $A(x+1,y+1)=A(x,A(x+1,y))$:
(a) calcula $A(1,2)$ y $A(2,1)$ aplicando \emph{solo} la definición;
(b) demuestra por inducción sobre $z$ que $A(1,z)=z+2$ y $A(2,z)=2z+3$.
\end{prob}
\begin{sol}
(a) $A(1,0)=A(0,1)=2$; \ $A(1,1)=A(0,A(1,0))=A(0,2)=3$; \ $A(1,2)=A(0,A(1,1))=A(0,3)=\mathbf 4$.\\
$A(2,0)=A(1,1)=3$; \ $A(2,1)=A(1,A(2,0))=A(1,3)=A(0,A(1,2))=A(0,4)=\mathbf 5$.

(b) \textbf{$A(1,z)=z+2$.} Base: $A(1,0)=A(0,1)=2$. H.I.: $A(1,n)=n+2$. Paso:
$A(1,n+1)=A(0,A(1,n))=A(0,n+2)=n+3=(n+1)+2$. \checkmark

\textbf{$A(2,z)=2z+3$.} Base: $A(2,0)=A(1,1)=3$. H.I.: $A(2,n)=2n+3$. Paso:
$A(2,n+1)=A(1,A(2,n))=A(1,2n+3)=(2n+3)+2=2(n+1)+3$, usando el resultado anterior como lema. \checkmark
\end{sol}

\begin{prob}
Demuestra que $A(3,z)=2^{z+3}-3$. Calcula $A(3,3)$ y $A(4,1)$. ¿Por qué el programa recursivo de Ackermann termina en \texttt{StackOverflowException} para valores pequeños como $(4,2)$?
\end{prob}
\begin{sol}
Base: $A(3,0)=A(2,1)=5=2^3-3$. Paso: $A(3,n+1)=A(2,A(3,n))=2(2^{n+3}-3)+3=2^{n+4}-6+3=2^{(n+1)+3}-3$. \checkmark

$A(3,3)=2^6-3=\mathbf{61}$. \ $A(4,0)=A(3,1)=13$; \ $A(4,1)=A(3,A(4,0))=A(3,13)=2^{16}-3=\mathbf{65533}$.
En general $A(4,z)$ es una torre de $z+3$ doses menos 3; $A(4,2)=2^{65536}-3$ tiene 19\,729 dígitos.

La recursión es doble y el argumento interno crece de forma no primitiva recursiva: la profundidad de la pila y el número de llamadas son astronómicos antes de llegar a un caso base; ninguna pila finita alcanza.
\end{sol}

\begin{prob}
Usando los lemas del problema anterior, demuestra $\hoare{x=2}{w\asg A(x,y)+1}{w=2y+4}$ y $\hoare{x=1}{w\asg A(x,y)}{w=y+2}$. ¿Por qué no basta el axioma de asignación?
\end{prob}
\begin{sol}
Primera: por asignación, $\sust{(w=2y+4)}{A(x,y)+1}{w}\equiv A(x,y)+1=2y+4\equiv A(x,y)=2y+3$, así $\hoare{A(x,y)=2y+3}{w\asg A(x,y)+1}{w=2y+4}$. Por el lema $A(2,z)=2z+3$: $x=2\Rightarrow A(x,y)=2y+3$. Por consecuencia:
\[\dfrac{x=2\Rightarrow A(x,y)=2y+3\qquad\hoare{A(x,y)=2y+3}{w\asg A(x,y)+1}{w=2y+4}}{\hoare{x=2}{w\asg A(x,y)+1}{w=2y+4}}.\]
Segunda: idéntica con $\hoare{A(x,y)=y+2}{w\asg A(x,y)}{w=y+2}$ y el lema $x=1\Rightarrow A(x,y)=y+2$.

El axioma solo produce la precondición $A(x,y)=\dots$; para llegar a $x=1$ o $x=2$ hace falta un \emph{teorema} sobre $A$ (demostrado por inducción aparte) y la regla de consecuencia. Así se reutilizan lemas sin volver a desenrollar la recursión.
\end{sol}

\begin{prob}
Compara la inducción usada para ciclos con la usada para $f$ (McCarthy) y Ackermann: dirección, hipótesis y qué se demuestra.
\end{prob}
\begin{sol}
\begin{tabular}{@{}p{3cm}p{5.5cm}p{6.3cm}@{}}
\toprule & Ciclos & McCarthy / Ackermann\\\midrule
Dirección & sube desde $n=0$ (iteraciones) & McCarthy: baja desde un techo $k$; Ackermann: inducción anidada sobre $x$ y, dentro, sobre $y$\\
Hipótesis & $P(n)$ para el $n$ anterior & $P(m)$ para \emph{todo} $m\in(x,k]$ (fuerte)\\
Qué prueba & un invariante se conserva & la función devuelve cierto valor/fórmula\\
Ejemplo & $C=DA$ en \texttt{Cuadrado} & $f(x)=91$; $A(2,z)=2z+3$\\\bottomrule
\end{tabular}
\end{sol}

% =====================================================================
\section{Programación lógica y Prolog}
% =====================================================================

\begin{prob}
(a) Escribe como cláusula (disyunción de literales) y como implicación la regla \texttt{abuelo(X,Z) :- padre(X,Y), padre(Y,Z).} (b) ¿Qué es una cláusula de Horn definida? ¿Y una cláusula objetivo? Representa la consulta \texttt{?- abuelo(juan, W).}
\end{prob}
\begin{sol}
(a) Implicación: $\forall X,Y,Z\;(padre(X,Y)\land padre(Y,Z)\to abuelo(X,Z))$. Cláusula: $abuelo(X,Z)\lor\neg padre(X,Y)\lor\neg padre(Y,Z)$.

(b) \textbf{Horn}: a lo más un literal positivo. \textbf{Definida}: exactamente uno (hechos y reglas; la cabeza es el positivo). \textbf{Objetivo}: ningún literal positivo; la consulta es $\leftarrow abuelo(juan,W)\equiv\neg abuelo(juan,W)$. Un hecho $padre(juan,pedro)$ es $padre(juan,pedro)\leftarrow\T$.
\end{sol}

\begin{prob}
Da el unificador más general o explica por qué no existe:
(a) $p(X,b)$ y $p(a,Y)$; \ (b) $f(X,g(Y))$ y $f(h(Z),g(a))$; \ (c) $p(X,X)$ y $p(a,b)$; \ (d) $gusta(maria,X)$ y $gusta(Y,cine)$; \ (e) $q(f(X),X)$ y $q(Y,a)$; \ (f) $p(X)$ y $p(f(X))$.
\end{prob}
\begin{sol}
(a) $\{X/a,\ Y/b\}$. \ (b) $\{X/h(Z),\ Y/a\}$. \ (c) No: $X/a$ obliga a $a=b$, constantes distintas. \ (d) $\{Y/maria,\ X/cine\}$. \ (e) $\{Y/f(a),\ X/a\}$ (tras $X/a$, $f(X)$ se vuelve $f(a)$). \ (f) No (prueba de ocurrencia): $X=f(X)$ exigiría un término infinito.
\end{sol}

\begin{prob}
Demuestra por \textbf{refutación} con resolución:
(a) $S=\{\neg hombre(X)\lor mortal(X),\ hombre(socrates)\}\models mortal(socrates)$.
(b) $S=\{P(X)\lor Q(X),\ \neg Q(X)\lor R(X),\ \neg R(X)\lor P(X)\}\models P(X)$.
\end{prob}
\begin{sol}
Se agrega la negación de la conclusión y se busca la cláusula vacía $\square$.

(a) $C_1:\neg hombre(X)\lor mortal(X)$, $C_2: hombre(socrates)$, $C_3:\neg mortal(socrates)$.
$C_3,C_1$ con $\theta=\{X/socrates\}$: $\neg hombre(socrates)$; con $C_2$: $\square$. Luego $S\models mortal(socrates)$.

(b) $C_4:\neg P(X)$.
$C_4,C_3\Rightarrow\neg R(X)$; \ $\neg R(X),C_2\Rightarrow\neg Q(X)$; \ $\neg Q(X),C_1\Rightarrow P(X)$; \ $P(X),C_4\Rightarrow\square$.
$S\cup\{\neg P(X)\}$ es insatisfactible $\Rightarrow S\models P(X)$.
\end{sol}

\begin{prob}
Dado el programa
\begin{lstlisting}[language=Prolog]
padre(juan, pedro).
padre(pedro, luis).
padre(pedro, ana).
abuelo(X, Z) :- padre(X, Y), padre(Y, Z).
\end{lstlisting}
(a) Escribe la derivación SLD para \texttt{?- abuelo(juan, W).} indicando en cada paso la cláusula usada y el unificador. (b) ¿Qué responde Prolog al teclear \texttt{;}?
\end{prob}
\begin{sol}
Renombramos la regla: $abuelo(X_1,Z_1)\leftarrow padre(X_1,Y_1),padre(Y_1,Z_1)$.
\begin{enumerate}
\item $G_0:\ \leftarrow abuelo(juan,W)$. Con la regla, $\theta_1=\{X_1/juan,\ Z_1/W\}$.
\item $G_1:\ \leftarrow padre(juan,Y_1),\ padre(Y_1,W)$. Primera meta con el hecho $padre(juan,pedro)$: $\theta_2=\{Y_1/pedro\}$.
\item $G_2:\ \leftarrow padre(pedro,W)$. Con $padre(pedro,luis)$: $\theta_3=\{W/luis\}$.
\item $G_3:\ \square$. Respuesta: \texttt{W = luis}.
\end{enumerate}
(b) Con \texttt{;} Prolog retrocede al último punto de elección ($G_2$) y usa $padre(pedro,ana)$: \texttt{W = ana}. Con otro \texttt{;} no quedan cláusulas para $padre(pedro,W)$ ni para $padre(juan,Y_1)$: \texttt{false}.
\end{sol}

\begin{prob}
Explica por qué el paso ``reemplazar una meta por el cuerpo de una regla'' es una instancia de la regla de resolución. ¿Qué significan las letras S, L, D?
\end{prob}
\begin{sol}
La regla $A\leftarrow B_1,\dots,B_m$ es la cláusula $A\lor\neg B_1\lor\dots\lor\neg B_m$ y la meta $\leftarrow A_1,\dots,A_k,\dots,A_n$ es $\neg A_1\lor\dots\lor\neg A_k\lor\dots\lor\neg A_n$. Si $\theta=\mathrm{UMG}(A,A_k)$, los literales $A$ y $\neg A_k$ son complementarios: se eliminan y el resolvente es
$(\neg A_1\lor\dots\lor\neg A_{k-1}\lor\neg B_1\lor\dots\lor\neg B_m\lor\neg A_{k+1}\lor\dots\lor\neg A_n)\theta$, que en notación de metas es $\leftarrow(A_1,..,A_{k-1},B_1,..,B_m,A_{k+1},..,A_n)\theta$.

\textbf{S}: regla de \emph{selección} de la meta (Prolog: la de más a la izquierda). \textbf{L}: \emph{lineal} (cada resolvente se obtiene del anterior y una cláusula del programa). \textbf{D}: cláusulas \emph{definidas}. Prolog además prueba las cláusulas en orden de programa, en profundidad, con retroceso.
\end{sol}

\begin{prob}
Dado el programa, indica si cada consulta tiene éxito, falla o produce error, y justifica:
\begin{lstlisting}[language=Prolog]
a.
b.
c :- fail.
s :- a, \+ c.
t :- c ; b.
u :- a, c.
v :- \+ a ; c.
w :- b, ( \+ c, a ).
x :- y.
\end{lstlisting}
Consultas: \texttt{?- s.} \texttt{?- t.} \texttt{?- u.} \texttt{?- v.} \texttt{?- w.} \texttt{?- x.}
\end{prob}
\begin{sol}
\begin{tabular}{@{}llp{11cm}@{}}
\texttt{s} & \checkmark & \texttt{a} es hecho; \texttt{c} no se puede demostrar (su cuerpo es \texttt{fail}), así que \texttt{\textbackslash+ c} tiene éxito.\\
\texttt{t} & \checkmark & disyunción: \texttt{c} falla, se prueba la alternativa \texttt{b}, que es hecho.\\
\texttt{u} & \texttimes & \texttt{a} tiene éxito pero \texttt{c} falla.\\
\texttt{v} & \texttimes & \texttt{\textbackslash+ a} falla (porque \texttt{a} sí se demuestra) y \texttt{c} también falla.\\
\texttt{w} & \checkmark & \texttt{b}; luego \texttt{\textbackslash+ c} tiene éxito y \texttt{a} es hecho.\\
\texttt{x} & error & \texttt{y} no está definido: SWI-Prolog lanza \emph{Unknown procedure: y/0} (\texttt{existence\_error}); no es lo mismo que ``falso''.
\end{tabular}

Recordatorio: \texttt{\textbackslash+ G} es \emph{negación como falla}: tiene éxito si el programa \emph{no logra demostrar} \texttt{G}.
\end{sol}

\begin{prob}
Con el programa \texttt{MariaGusta.pl}
\begin{lstlisting}[language=Prolog]
gusta( maria, helado ).
gusta( maria, leer ).
gusta( maria, cine ).
dar( juan, beso, maria ).
recibe( X, beso ) :- dar( juan, beso, X ).
\end{lstlisting}
indica la respuesta de: (a) \texttt{gusta(maria, cine).} (b) \texttt{gusta(juan, cine).} (c) \texttt{gusta(Quien, helado).} (d) \texttt{gusta(maria, Maria).} (con \texttt{;}) (e) \texttt{recibe(X, beso).} (f) \texttt{recibe ( X, beso ).}
\end{prob}
\begin{sol}
(a) \texttt{true} (átomos que coinciden). \ (b) \texttt{false}. \ (c) \texttt{Quien = maria}. \
(d) \texttt{Maria} es una \emph{variable} (mayúscula inicial), no el átomo \texttt{maria}: responde \texttt{Maria = helado ; Maria = leer ; Maria = cine}. \
(e) Se resuelve con la regla: la meta \texttt{dar(juan,beso,X)} unifica con el hecho: \texttt{X = maria}. \
(f) \textbf{Error de sintaxis} (\emph{operator expected}): no puede haber espacio entre el nombre del predicado y el paréntesis.
\end{sol}

\begin{prob}[\simil]
Considera
\begin{lstlisting}[language=Prolog]
maximo( X, Y, X ) :- X >= Y.
maximo( X, Y, Y ) :- X < Y.
mayor( X, Y, X ) :- X >= Y, !.
mayor( X, Y, Y ) :- X < Y.
max( X, Y, Y ) :- X < Y, !.
max( X, _, X ).
\end{lstlisting}
(a) ¿Qué hace el corte \texttt{!}? (b) ¿Qué es un corte ``verde'' y uno ``rojo''? Clasifica los de \texttt{mayor} y \texttt{max}. (c) ¿Qué responde \texttt{?- max(3, 5, 3).}? ¿Es correcto? Corrígelo.
\end{prob}
\begin{sol}
(a) El corte tiene éxito una vez y \emph{elimina los puntos de elección} creados desde que se eligió la cláusula actual: Prolog ya no probará otras cláusulas del mismo predicado ni otras alternativas de las metas anteriores al \texttt{!} dentro del cuerpo.

(b) \textbf{Verde:} solo mejora eficiencia; quitarlo no cambia las respuestas. \textbf{Rojo:} cambia el significado; quitarlo cambia las respuestas. En \texttt{mayor} es verde (las condiciones $X\ge Y$ y $X<Y$ son excluyentes; el corte evita probar una cláusula que igual fallaría). En \texttt{max} es rojo: la segunda cláusula no tiene condición y depende del corte (sin él, \texttt{max(3,5,M)} daría \texttt{M = 5 ; M = 3}).

(c) \texttt{?- max(3,5,3).} La primera cláusula exige unificar el tercer argumento con $Y=5$: $3\ne5$, falla \emph{antes} de llegar al corte. La segunda, \texttt{max(X,\_,X)}, unifica con $X=3$: responde \texttt{true}. \textbf{Incorrecto} (el máximo es 5).
Corrección: hacer la unificación de la salida \emph{después} del corte:
\begin{lstlisting}[language=Prolog]
max( X, Y, Z ) :- X < Y, !, Z = Y.
max( X, _, X ).
\end{lstlisting}
o usar la versión \texttt{maximo} sin cortes.
\end{sol}

\begin{prob}
Indica la respuesta de SWI-Prolog: (a) \texttt{X = 2+3.} (b) \texttt{X is 2+3.} (c) \texttt{5 is 2+3.} (d) \texttt{2+3 = 5.} (e) \texttt{2+3 =:= 5.} (f) \texttt{X is Y+1.}
\end{prob}
\begin{sol}
(a) \texttt{X = 2+3} (unifica con el \emph{término} sin evaluar). (b) \texttt{X = 5} (\texttt{is} evalúa el lado derecho). (c) \texttt{true}. (d) \texttt{false} (el término \texttt{+(2,3)} no es el número 5). (e) \texttt{true} (\texttt{=:=} compara valores aritméticos). (f) Error: \emph{arguments are not sufficiently instantiated} (\texttt{Y} no tiene valor).
\end{sol}

\begin{prob}
(a) Traza \texttt{?- factorial(3, F).} con
\begin{lstlisting}[language=Prolog]
factorial(0, 1) :- !.
factorial(N, F) :- N > 0, N1 is N - 1, factorial(N1, F1), F is N * F1.
\end{lstlisting}
(b) Escribe \texttt{suma\_hasta(N, S)} ($S=1+2+\dots+N$) y \texttt{mcd(A, B, G)} (algoritmo de Euclides) con el mismo estilo.
\end{prob}
\begin{sol}
(a) \texttt{factorial(3,F)}: $3>0$, $N1=2$, llama \texttt{factorial(2,F1)} $\to$ $N1=1$, \texttt{factorial(1,F1')} $\to$ $N1=0$, \texttt{factorial(0,F1'')}: primera cláusula, $F1''=1$ y el corte descarta la segunda. Al regresar: $F1'=1\cdot1=1$, $F1=2\cdot1=2$, $F=3\cdot2=\mathbf 6$.
\begin{lstlisting}[language=Prolog]
suma_hasta(0, 0) :- !.
suma_hasta(N, S) :- N > 0, N1 is N - 1, suma_hasta(N1, S1), S is S1 + N.

mcd(A, 0, A) :- !.
mcd(A, B, G) :- B > 0, R is A mod B, mcd(B, R, G).
\end{lstlisting}
\texttt{?- mcd(48, 18, G).} $\to$ \texttt{mcd(18,12)} $\to$ \texttt{mcd(12,6)} $\to$ \texttt{mcd(6,0)} $\to$ \texttt{G = 6}. Mismo patrón que en clase: caso base protegido por corte y caso recursivo que reduce el problema.
\end{sol}

\begin{prob}
Compara resolución general y resolución SLD (qué cláusulas usan, forma de la derivación, qué significa el éxito, cómo manejan alternativas). ¿Por qué el orden de las cláusulas importa en Prolog aunque lógicamente no importe?
\end{prob}
\begin{sol}
\begin{tabular}{@{}p{3.3cm}p{5.8cm}p{6cm}@{}}
\toprule & General & SLD / Prolog\\\midrule
Cláusulas & arbitrarias & definidas $+$ una meta\\
Elección & cualquier par de cláusulas & la meta seleccionada con una cláusula del programa\\
Derivación & árbol/grafo amplio & lineal por rama\\
Éxito & se deriva $\square$ & la lista de metas queda vacía\\
Alternativas & muchas combinaciones & retroceso (backtracking)\\\bottomrule
\end{tabular}

Lógicamente el conjunto de consecuencias no cambia al reordenar. Operacionalmente Prolog elige la meta de más a la izquierda y prueba cláusulas en orden, en profundidad: un orden desafortunado puede ser lento o \emph{no terminar} (p.\,ej.\ una regla recursiva por la izquierda como \texttt{anc(X,Y) :- anc(X,Z), padre(Z,Y).} antes del caso base). Además, que una búsqueda no termine no significa que la consulta sea falsa.
\end{sol}
